\documentclass[a4paper, 10pt]{article}
\usepackage[utf8]{inputenc} % Change according your file encoding
\usepackage{graphicx}
\usepackage{url}

\usepackage{hyperref}
\hypersetup{
    linktocpage,
    pdfborderstyle={/S/U/W 1},
    pdfborder=0 0 1,
}

\usepackage[nameinlink, noabbrev]{cleveref}
\usepackage{booktabs}
\usepackage{caption}
\usepackage{subcaption}
\captionsetup{labelfont=sc,justification=centerlast}
\captionsetup[table]{skip=3pt}
\subcaptionsetup{labelfont=up}

\usepackage{microtype}

\usepackage{titlesec}
% \titleformat*{\subsection}{\itshape}
% \renewcommand\thesubsection{\roman{subsection})}

\usepackage[newfloat]{minted2}
\setminted{
    breaklines,
    fontfamily=tt,
    style=bw,
    escapeinside=¬¬,
    tabsize=4,
}
\setmintedinline{
    fontfamily=tt,
    style=bw,
}

\RequirePackage{xcolor}
\definecolor{gris}{HTML}{747474}
\definecolor{azul}{HTML}{19327F}
\definecolor{rojo}{HTML}{BE511C}
\definecolor{turquesa}{HTML}{0F7699}
\definecolor{purpura}{HTML}{8B2368}
\definecolor{morado}{HTML}{2C1D41}
\definecolor{verde}{HTML}{4B7463}

\usepackage{pgfplots}

\pgfplotsset{
    % Ejes
    axis x line=bottom,
    axis y line=left,
    axis line style={shorten >= -6pt}, % Alargar flechas
    /pgfplots/layers/custom/.define layer set={
        axis background, axis grid, main, axis ticks, axis lines, axis tick labels, axis descriptions, axis foreground
    }{/pgfplots/layers/standard},
    set layers=custom,
    % Ticks
    minor x tick num=1,
    minor y tick num=1,
    every x tick/.style={color=black, thin},
    every y tick/.style={color=black, thin},
    every minor x tick/.style={color=gris, thin},
    every minor y tick/.style={color=gris, thin},
    grid style={color=gris, loosely dashed, line cap=rect},
    % Estilo
    every axis plot/.append style={
        very thick,
        line cap=rect,
    },
    % Leyenda
    legend cell align=left,
    %legend style={draw=none},
    % Separador de miles
    /pgf/number format/.cd,
    1000 sep={\,},
}
\pgfplotsset{
    cycle multi list={
        {solid, mark=*}, \nextlist azul, rojo, turquesa, purpura, morado
    }
}


%opening
\title{Seminar Report: Opty}
\author{\textbf{Castilla Rodríguez, Víctor}\and \textbf{Chavez Quevedo, Luis Fernando}\and \textbf{Amorós Cubells, Alfonso}
}
\date{\normalsize\today{}}




\begin{document}

\maketitle

% \begin{center}
%   Upload your report in PDF format.

%   Use this LaTeX template to format the report.

% 	A compressed file (.tar.gz) containing all your source code files must be submitted together with this report\footnote{Describe in the report any design decision required to understand your code (if any)}.
% \end{center}

\section{Introduction}

% \textit{Introduce in a couple of sentences the seminar and the main topic related to distributed systems it covers.}


The objective of this seminar is implementing  optimistic concurrency control in Erlang, experimenting with different parameters (number of clients, store size, read/writes per transaction...) and seeing how they affect transactions.

The provided files contain three folders: \texttt{opty}, for the main implementation of Opty; \texttt{opty-subsets}, for the Opty implementation of \cref{sec:ptge}; and \texttt{timey-subsets}, for the Timey implementation of \cref{sec:timey-subsets}.

Along with the code files, an auxiliary Makefile is provided for building and
running the experiments. Thus, one can call \texttt{make run} to run the code.

\section{Performance} \label{sec:performance}

% \textit{Provide evidence of the experiments you did (e.g., use screenshots) and discuss the results you got. In addition, you may provide figures or tables with experimental results of the system evaluation. For each seminar, we will provide you with some guidance on which kind of evaluation you should do.}

The following subsections explore how varying parameters affects the percentage of transactions that commit successfully.
Unless otherwise specified, the default parameters are 5 clients, 50 entries, 1 read and 1 write per transaction, and 5 seconds of execution time.

The percentage of successful transactions shown is calculated as the average of the success rate of every client.
It is important to note that, except for \cref{sec:ptge}, the success rate is almost the same for all clients in every experiment (standard deviation is \(< 0.5\%\)).

\subsection{Number of concurrent clients}

\Cref{fig:nclients} shows the correlation between the number of clients and the average successful transactions.
With 1 client, 100\% of operations succeed, as there can never be conflicts.
This percentage steadily decreases with the increase of concurrent clients.

With 1 read and 1 write per transaction (labeled in the graph as 1\,RD 1\,WR), it falls down to 87.8\% with 10 clients and 68.8\% with 30 clients.

With 2 reads and 2 writes,
conflicts are more likely to appear, causing transactions to abort more often.
The success rate is 68.1\% with 10 clients and 41.6\% with 30 clients.

\begin{figure}[htp]
\begin{center}
\begin{tikzpicture}
\begin{axis}[
    ymajorgrids=true,
    xmin=0,
    xmax=30,
    ymin=0,
    ymax=100,
    xlabel=Number of clients,
    ylabel=Successful transactions (\%),
    legend pos=south east,
    width=.98\textwidth,
    height=0.45\textwidth,
]

\addplot table [
    col sep=semicolon,
    x=Clients,
    y=Avg,
    /pgf/number format/read comma as period,
] {csv/clients_50_1rd1wr.csv};
\addlegendentry{1\,RD, 1\,WR}

\addplot table [
    col sep=semicolon,
    x=Clients,
    y=Avg,
    /pgf/number format/read comma as period,
] {csv/clients_50_2rd2wr.csv};
\addlegendentry{2\,RD, 2\,WR}

\end{axis}
\end{tikzpicture}
\caption{Average successful transactions depending on the  number of clients.}
\label{fig:nclients}
\end{center}
\end{figure}


\subsection{Number of entries}

% Baseline: 5 clients 1RD 1WR
\Cref{fig:nentries} shows how the number of entries in the store impacts successful transactions.
As one might expect in this scenario, increasing the number of entries reduces the possibility of two transactions picking the same entry and thus increases the average number of successful transactions.

The graph shows this is true both for 1\,RD 1\,WR per transaction and for 2\,RD 2\,WR per transaction (although there is an ``edge case'' with just 1 entry, in which both cases are conceptually identical but with different overhead).
With 50 entries (not shown in the graph), the averages are 95\% and 84\%, respectively.

\begin{figure}[htp]
\begin{center}
\begin{tikzpicture}
\begin{axis}[
    ymajorgrids=true,
    xmin=0,
    xmax=30.2,
    ymin=0,
    ymax=100,
    xlabel=Number of entries,
    ylabel=Successful transactions (\%),
    legend pos=south east,
    width=.98\textwidth,
    height=0.45\textwidth,
]

\addplot table [
    col sep=semicolon,
    x=Entries,
    y=Avg,
    /pgf/number format/read comma as period,
] {csv/nentries (copia).csv};
\addlegendentry{1\,RD, 1\,WR}

\addplot table [
    col sep=semicolon,
    x=Entries,
    y=Avg,
    /pgf/number format/read comma as period,
] {csv/nentries_2RD2WR (copia).csv};
\addlegendentry{2\,RD, 2\,WR}

\end{axis}
\end{tikzpicture}
\caption{Average successful transactions depending on the  number of entries.}
\label{fig:nentries}
\end{center}
\end{figure}


\subsection{Number of read operations per transaction}

\Cref{fig:nreads} shows the effect of varying the number of reads per transaction (with 1 write).
There is a linear decrease in successful transactions as the number of reads increases, falling down from 100\% with 0 reads to 60.3\% with 20 reads.

\begin{figure}[htp]
\begin{center}
\begin{tikzpicture}
\begin{axis}[
    ymajorgrids=true,
    xmin=0,
    % xmax=30.2,
    ymin=0,
    ymax=100,
    xlabel=Number of reads,
    ylabel=Successful transactions (\%),
    legend pos=south east,
    width=.98\textwidth,
    height=0.45\textwidth,
]

\addplot table [
    col sep=semicolon,
    x=Rd,
    y=Avg,
    /pgf/number format/read comma as period,
] {csv/reads.csv};
% \addlegendentry{1\,RD, 1\,WR}

\end{axis}
\end{tikzpicture}
\caption{Average successful transactions depending on the  number of reads.}
\label{fig:nreads}
\end{center}
\end{figure}


\subsection{Number of write operations per transaction}

\Cref{fig:nwrites} shows the effect of varying the number of writes per transaction (with 1 read).
Note that increasing the number of writes makes performance worsen faster than in the previous scenario of increasing the number of reads.
That is, the performance of 1\,RD 20\,WR is worse than 20\,RD 1\,WR.

\begin{figure}[htp]
\begin{center}
\begin{tikzpicture}
\begin{axis}[
    ymajorgrids=true,
    xmin=0,
    % xmax=30.2,
    ymin=0,
    ymax=100,
    xlabel=Number of writes,
    ylabel=Successful transactions (\%),
    legend pos=south east,
    width=.98\textwidth,
    height=0.45\textwidth,
]

\addplot table [
    col sep=semicolon,
    x=Wr,
    y=Avg,
    /pgf/number format/read comma as period,
] {csv/writes.csv};
% \addlegendentry{1\,RD, 1\,WR}

\end{axis}
\end{tikzpicture}
\caption{Average successful transactions depending on the  number of writes.}
\label{fig:nwrites}
\end{center}
\end{figure}

\subsection{Ratio of read and write operations per transaction}

\Cref{fig:ratio} displays the performance of varying the ratio of read and write operations
for a fixed amount of 10 operations per transaction (i.e., 10\,RD 0\,WR at the beginning and 0\,RD 10\,WR at the end).

As expected, the extreme cases have 100\% successful transactions, as no conflicts can occur.
The worst performance happens exactly at the midway point (55.12\%), although 4\,RD 6\,WR has similar bad performance (55.17\%).
Consistent with the previous two subsections, writes are worse for performance, so for instance 1\,RD 9\,WR performs worse than 9\,RD 1\,WR.

\begin{figure}[htp]
\begin{center}
\begin{tikzpicture}
\begin{axis}[
    ymajorgrids=true,
    xmin=0,
    % xmax=30.2,
    ymin=0,
    ymax=100,
    xlabel=Ratio of write operations (\%),
    ylabel=Successful transactions (\%),
    legend pos=south east,
    width=.98\textwidth,
    height=0.45\textwidth,
]

\addplot table [
    col sep=semicolon,
    x=WrPercent,
    y=Avg,
    /pgf/number format/read comma as period,
] {csv/ratio10.csv};
% \addlegendentry{1\,RD, 1\,WR}

\end{axis}
\end{tikzpicture}
\caption{Performance under a different ratio of read/writes with 10 total operations.}
\label{fig:ratio}
\end{center}
\end{figure}


\subsection{Percentage of accessed entries} \label{sec:ptge}

To vary the percentage of accessed entries with
respect to the total number of entries,
we need to generate a random subset of entries for each client.
% (i.e. each client accesses a randomly
% generated subset of the store: the number of entries in each subset should be the
% same for all the clients, but the subsets should be different per client and should
% not contain necessarily contiguous entries)
To do this, the clients now take the \texttt{SubsetLen} parameter (instead of the number of \texttt{Reads} and \texttt{Writes}) to generate a random subset:
\begin{minted}{erlang}
start(ClientID, Entries, Server, SubsetLen) ->
    Subset = lists:sublist([X || {_, X} <- lists:sort([{rand:uniform(), E} || E <- lists:seq(1, Entries)])], SubsetLen),
    spawn(fun() -> open(ClientID, Subset, Server, 0, 0) end).
\end{minted}

Then, when performing a transaction, only entries from the subset are read or written:
\begin{minted}{erlang}
do_transaction(_, [], Handler) ->
    do_commit(Handler);
do_transaction(ClientID, Subset, Handler) ->
    [Num | RemSubset] = Subset,
    Op = rand:uniform(),
    if Op >= 0.5 ->
         do_read(Num, Handler);
       true ->
         do_write(Num, Handler, ClientID)
    end,
    do_transaction(ClientID, RemSubset, Handler).
\end{minted}

\Cref{fig:vi} shows the effect of varying the subset size up to 30 (i.e., 60\% of entries).
Unlike previous experiments, the resulting success rate varies greatly across clients: some clients might never experience conflicts if all of their entries are different from other clients, while others might share a lot of entries and thus have many conflicts.
This is shown in the graph by displaying the maximum and minimum success rates obtained as error bars.

\begin{figure}[htp]
\begin{center}
\begin{tikzpicture}
\begin{axis}[
    ymajorgrids=true,
    xmin=0,
    xmax=30.2,
    ymin=0,
    ymax=100,
    xlabel=Subset size,
    ylabel=Successful transactions (\%),
    legend pos=south east,
    width=.98\textwidth,
    height=0.46\textwidth,
]

\addplot plot[error bars/.cd, y dir=both, y explicit] table [
    col sep=semicolon,
    x=SubsetSize,
    y=Avg,
    y error plus=MaxMinusAvg,
    y error minus=AvgMinusMin,
    /pgf/number format/read comma as period,
] {csv/subsetCopy.csv};
% \addlegendentry{1\,RD, 1\,WR}

\end{axis}
\end{tikzpicture}
\caption{Performance while varying the subset size for every client.}
\label{fig:vi}
\end{center}
\end{figure}

The number of conflicts depends greatly on the random subsets that have been chosen.
For instance, with subset size 4, one client had a success rate of 100\% while the average was 76.6\%.
Another extreme example ocurred with subset size 8, where one client had a success rate of 87\% while another only had 35\% (the average was 65.5\%). % with a standard deviation of 20.46

As the subset size increases, ``unique'' sets for clients are less likely, and thus the percentage of successful transactions is roughly the same for every client.
This is specially evident if we look at the standard deviation of the success percentage, shown in \cref{fig:stdv}.

\begin{figure}[htp]
\begin{center}
\begin{tikzpicture}
\begin{axis}[
    ymajorgrids=true,
    xmin=0,
    xmax=30.2,
    ymin=0,
    ymax=22,
    xlabel=Subset size,
    ylabel=Standard deviation (\%),
    legend pos=south east,
    width=.98\textwidth,
    height=0.4\textwidth,
]

\addplot plot[error bars/.cd, y dir=both, y explicit] table [
    col sep=semicolon,
    x=SubsetSize,
    y=DesvTip,
    % y error plus=MaxMinusAvg,
    % y error minus=AvgMinusMin,
    /pgf/number format/read comma as period,
] {csv/subsetCopy.csv};
% \addlegendentry{1\,RD, 1\,WR}

\end{axis}
\end{tikzpicture}
\caption{Standard deviation of successful transactions depending on subset size.}
\label{fig:stdv}
\end{center}
\end{figure}

\section{Other concurrency control techniques:\\ Timestamp ordering}

In this section, the previously shown performance of the optimistic concurrency control method will be compared and evaluated with the timestamp ordering technique's performance. As such, the provided implementation of timestamp ordering (Timey) will be used against the optimistic concurrency control's implementation (Opty).

As with \cref{sec:performance}, the default parameters are 5 clients, 50 entries, 1 read and 1 write per transaction, and 5 seconds of execution time.

\subsection{Number of clients}

\Cref{fig:tnclients} shows that the timestamp ordering technique averages a higher number of successful transactions than the optimistic method does.

With 1 read and 1 write, the average success of timestamp ordering falls down to 84.5\% when working with 30 concurrent clients. Meanwhile, with 2 reads and 2 writes, the performance of timestamp ordering averages the same number of successful transactions as the optimistic method with 1 read and 1 write.

\begin{figure}[htp]
\begin{center}
\begin{tikzpicture}
\begin{axis}[
    ymajorgrids=true,
    xmin=0,
    xmax=30,
    ymin=0,
    ymax=100,
    xlabel=Number of clients,
    ylabel=Successful transactions (\%),
    legend pos=south west,
    width=.98\textwidth,
    height=0.45\textwidth,
]

\addplot table [
    col sep=semicolon,
    x=Clients,
    y=Avg,
    /pgf/number format/read comma as period,
] {csv/clients_50_1rd1wr.csv};
\addlegendentry{Opty 1\,RD, 1\,WR}

\addplot table [
    col sep=semicolon,
    x=Clients,
    y=Avg,
    /pgf/number format/read comma as period,
] {csv/clients_50_2rd2wr.csv};
\addlegendentry{Opty 2\,RD, 2\,WR}

\addplot table [
    col sep=semicolon,
    x=Clients,
    y=Avg,
    /pgf/number format/read comma as period,
] {csv/timey_clients_50_1rd1wr.csv};
\addlegendentry{Timey 1\,RD, 1\,WR}

\addplot table [
    col sep=semicolon,
    x=Clients,
    y=Avg,
    /pgf/number format/read comma as period,
] {csv/timey_clients_50_2rd2wr.csv};
\addlegendentry{Timey 2\,RD, 2\,WR}

\end{axis}
\end{tikzpicture}
\caption{Average successful transactions depending on the number of clients with optimistic concurrency control (Opty) and timestamp ordering (Timey).}
\label{fig:tnclients}
\end{center}
\end{figure}

\subsection{Number of entries}

\Cref{fig:tnentries} shows similar differences in transaction success rate between both concurrency control techniques as displayed in \cref{fig:tnclients}. Timestamp ordering delivers a higher number of successful transactions than the optimistic method. Even while benefiting from a higher number of entries, and thus, less conflicts, timestamp ordering delivers a higher performance.

With 1 read and 1 write, timestamp ordering delivers a higher number of successful transactions. And as shown in the previous section, with 2 reads and 2 writes, its performance drops down to the same level as an optimistic concurrency control with only 1 read and 1 write per transaction.

\begin{figure}[htp]
\begin{center}
\begin{tikzpicture}
\begin{axis}[
    ymajorgrids=true,
    xmin=0,
    xmax=30.2,
    ymin=0,
    ymax=100,
    xlabel=Number of entries,
    ylabel=Successful transactions (\%),
    legend pos=south east,
    width=.98\textwidth,
    height=0.45\textwidth,
]

\addplot table [
    col sep=semicolon,
    x=Entries,
    y=Avg,
    /pgf/number format/read comma as period,
] {csv/nentries (copia).csv};
\addlegendentry{Opty 1\,RD, 1\,WR}

\addplot table [
    col sep=semicolon,
    x=Entries,
    y=Avg,
    /pgf/number format/read comma as period,
] {csv/nentries_2RD2WR (copia).csv};
\addlegendentry{Opty 2\,RD, 2\,WR}

\addplot table [
    col sep=semicolon,
    x=Entries,
    y=Avg,
    /pgf/number format/read comma as period,
] {csv/timey_nentries_1rd1wr.csv};
\addlegendentry{Timey 1\,RD, 1\,WR}

\addplot table [
    col sep=semicolon,
    x=Entries,
    y=Avg,
    /pgf/number format/read comma as period,
] {csv/timey_nentries_2rd2wr.csv};
\addlegendentry{Timey 2\,RD, 2\,WR}

\end{axis}
\end{tikzpicture}
\caption{Average successful transactions depending on the  number of entries with optimistic concurrency control (Opty) and timestamp ordering (Timey).}
\label{fig:tnentries}
\end{center}
\end{figure}

\subsection{Number of read operations per transaction}

For the data shown in this subsection, the number of write operations per transaction has been set to 1.

\Cref{fig:tnreads} reveals that the timestamp ordering implementation is able to maintain a higher average of successful transactions than the optimistic im\-ple\-men\-ta\-tion, even with a high number of read operations, dropping its success rate from 99.8\% with 0 read operations to 93.8\% with 20 read operations per transaction. Meanwhile, the optimistic implementation's performance drops down to 60.3\% when doing up to 20 read operations per transaction.

\begin{figure}[htp]
\begin{center}
\begin{tikzpicture}
\begin{axis}[
    ymajorgrids=true,
    xmin=0,
    % xmax=30.2,
    ymin=0,
    ymax=100,
    xlabel=Number of reads,
    ylabel=Successful transactions (\%),
    legend pos=south east,
    width=.98\textwidth,
    height=0.45\textwidth,
]

\addplot table [
    col sep=semicolon,
    x=Rd,
    y=Avg,
    /pgf/number format/read comma as period,
] {csv/reads.csv};
\addlegendentry{Opty}

\addplot table [
    col sep=semicolon,
    x=Rd,
    y=Avg,
    /pgf/number format/read comma as period,
] {csv/timey_reads.csv};
\addlegendentry{Timey}

\end{axis}
\end{tikzpicture}
\caption{Average successful transactions depending on the  number of reads with optimistic concurrency control (Opty) and timestamp ordering (Timey).}
\label{fig:tnreads}
\end{center}
\end{figure}

\subsection{Number of write operations per transaction}

\Cref{fig:tnwrites} show that the performance of both techniques, timestamp ordering and optimistic control, perform similary with different amounts of write operations per transaction.

\begin{figure}[htp]
\begin{center}
\begin{tikzpicture}
\begin{axis}[
    ymajorgrids=true,
    xmin=0,
    % xmax=30.2,
    ymin=0,
    ymax=100,
    xlabel=Number of writes,
    ylabel=Successful transactions (\%),
    legend pos=south east,
    width=.98\textwidth,
    height=0.45\textwidth,
]

\addplot table [
    col sep=semicolon,
    x=Wr,
    y=Avg,
    /pgf/number format/read comma as period,
] {csv/writes.csv};
\addlegendentry{Opty}

\addplot table [
    col sep=semicolon,
    x=Wr,
    y=Avg,
    /pgf/number format/read comma as period,
] {csv/timey_writes.csv};
\addlegendentry{Timey}

\end{axis}
\end{tikzpicture}
\caption{Average successful transactions depending on the  number of writes with optimistic concurrency control (Opty) and timestamp ordering (Timey).}
\label{fig:tnwrites}
\end{center}
\end{figure}

\subsection{Ratio of read and write operations per transaction}

\Cref{fig:tratio} displays the different possible combinations of read and write operations per transaction, where each transaction does a total of 10 operations.

As it was shown in previous sections, timestamp ordering performs better than optimistic concurrency control at handling multiple concurrent read operations per transaction, while showing that both methods have similar performance when doing multiple write operations per transaction.

This difference in performance between read and write operations is shown in \cref{fig:tratio}, where a high ratio of read operations results in a higher success rate by the timestamp ordering implementation, while the higher ratio of write operations per transactions leaves similar levels of success rate for both techniques.

\begin{figure}[htp]
\begin{center}
\begin{tikzpicture}
\begin{axis}[
    ymajorgrids=true,
    xmin=0,
    % xmax=30.2,
    ymin=0,
    ymax=100,
    xlabel=Ratio of write operations (\%),
    ylabel=Successful transactions (\%),
    legend pos=south east,
    width=.98\textwidth,
    height=0.45\textwidth,
]

\addplot table [
    col sep=semicolon,
    x=WrPercent,
    y=Avg,
    /pgf/number format/read comma as period,
] {csv/ratio10.csv};
\addlegendentry{Opty}

\addplot table [
    col sep=semicolon,
    x=WrPercent,
    y=Avg,
    /pgf/number format/read comma as period,
] {csv/timey_ratio10.csv};
\addlegendentry{Timey}

\end{axis}
\end{tikzpicture}
\caption{Performance under a different ratio of read/writes with 10 total operations with optimistic concurrency control (Opty) and timestamp ordering (Timey).}
\label{fig:tratio}
\end{center}
\end{figure}

\subsection{Percentage of accessed entries} \label{sec:timey-subsets}

\Cref{fig:tvi} shows the effect of varying the entry subset size given to each client in both implementations: timestamp ordering and optimistic control. The timestamp ordering implementation is able to maintain a higher success rate than the optimistic implementation, up to around a subset size of 20 entries. After that, its performance drops below the optimistic concurrency control method.

Timestamp ordering also shows a high variability in the success rate of its clients, like the optimistic method. However, unlike the optimistic control, timestamp ordering keeps a higher degree of variability in the number of successful transactions of each client even with a larger subset size, as shown in \cref{fig:tstdv}.

\begin{figure}[htp]
\begin{center}
\begin{tikzpicture}
\begin{axis}[
    ymajorgrids=true,
    xmin=0,
    xmax=30.2,
    ymin=0,
    ymax=100,
    xlabel=Subset size,
    ylabel=Successful transactions (\%),
    legend pos=north east,
    width=.98\textwidth,
    height=0.46\textwidth,
]

\addplot plot[error bars/.cd, y dir=both, y explicit] table [
    col sep=semicolon,
    x=SubsetSize,
    y=Avg,
    y error plus=MaxMinusAvg,
    y error minus=AvgMinusMin,
    /pgf/number format/read comma as period,
] {csv/subsetCopy.csv};
\addlegendentry{Opty}

\addplot plot[error bars/.cd, y dir=both, y explicit] table [
    col sep=semicolon,
    x=SubsetSize,
    y=Avg,
    y error plus=MaxMinusAvg,
    y error minus=AvgMinusMin,
    /pgf/number format/read comma as period,
] {csv/timey_subsets.csv};
\addlegendentry{Timey}

\end{axis}
\end{tikzpicture}
\caption{Performance while varying the subset size for every client with optimistic concurrency control (Opty) and timestamp ordering (Timey).}
\label{fig:tvi}
\end{center}
\end{figure}

\begin{figure}[htp]
\begin{center}
\begin{tikzpicture}
\begin{axis}[
    ymajorgrids=true,
    xmin=0,
    xmax=30.2,
    ymin=0,
    ymax=22,
    xlabel=Subset size,
    ylabel=Standard deviation (\%),
    legend pos=north east,
    width=.98\textwidth,
    height=0.4\textwidth,
]

\addplot plot[error bars/.cd, y dir=both, y explicit] table [
    col sep=semicolon,
    x=SubsetSize,
    y=DesvTip,
    % y error plus=MaxMinusAvg,
    % y error minus=AvgMinusMin,
    /pgf/number format/read comma as period,
] {csv/subsetCopy.csv};
\addlegendentry{Opty}

\addplot plot[error bars/.cd, y dir=both, y explicit] table [
    col sep=semicolon,
    x=SubsetSize,
    y=DesvTip,
    % y error plus=MaxMinusAvg,
    % y error minus=AvgMinusMin,
    /pgf/number format/read comma as period,
] {csv/timey_subsets.csv};
\addlegendentry{Timey}

\end{axis}
\end{tikzpicture}
\caption{Standard deviation of successful transactions depending on subset size with optimistic concurrency control (Opty) and timestamp ordering (Timey).}
\label{fig:tstdv}
\end{center}
\end{figure}

\section{Personal opinion}

% \textit{Provide your personal opinion of the seminar, indicating whether it should be included in next year's course or not.}

This seminar has proven to be a valuable learning experience for understanding optimistic concurrency control and timestamp ordering algorithms.
The provided Erlang code made the implementation straightforward and provided a simple way to experiment by tweaking each parameter individually. 
Overall, the seminar was useful and we believe it should be included in next year's course.

\end{document}
